Elastic LLM-serving policies depend on when capacity becomes available, not just
on request service times. We present and implement FidelityLoop, a closed-loop
validation framework connecting replay predictions, physical lifecycle execution,
and deadline-aware resource accounting. In a preregistered study of 36
Qwen2.5-7B-Instruct runs on two A100 GPUs, the calibrated replay improves
matched-request timing while worsening occupied-resource prediction; trajectory
analysis explains why aggregate error is insufficient for closed-loop evaluation.
We then train PPO scheduling policies in replay under two startup-time
assumptions and freeze five checkpoints for a 32-run deployment campaign (E2).
Three checkpoints meet every deadline in all six runs; two fail all six.
Both training environments produce passing and failing checkpoints.
Replay with calibrated service times and the historical startup proxy (C)
predicts feasibility correctly on 32/32 runs, versus 20/32 after changing
startup to 30 seconds (C30). Their cost MAEs are 0.0804 and 0.4657 scenario USD.
A follow-up on six new checkpoints instead finds all six feasible, including
three rejected by C. These cases show why both policies and their replay
models need deployment validation: aggregate resource error and past feasibility
agreement can each obscure wrong decisions. FidelityLoop exposes these failures
under a shared policy contract, retaining negative results within the studied
workload, hardware, and synthetic-API setting.
